\documentclass[../../main.tex]{subfiles}

\begin{document}

\section*{附加习题：良序}\label{section:ch01-附加习题良序}
\addcontentsline{toc}{section}{附加习题：良序}

在以下习题中要求证明\cref{axiom:ch01-选择公理}、良序原理与极大原理的等价性。我们要特别指出的是，只在习题 7 中用到了\cref{axiom:ch01-选择公理}。

\begin{exercise}\label{exercise:ch01-附加良序-1}
定理（广义归纳定义原理）设 $J$ 是一个良序集，$C$ 是一个集合。以 $\mathcal{F}$ 表示将 $J$ 的所有截映入 $C$ 的函数的全体。对于给定的一个函数 $\rho: \mathcal{F} \longrightarrow C$，存在唯一一个函数 $h: J \longrightarrow C$，使得对于每一个 $\alpha \in J$，$h(\alpha) = \rho(h \mid S_{\alpha})$。

[提示：仿照第\cref{section:ch01-良序集}节习题 10 中所给的证明思路。]
\end{exercise}

\begin{exercise}\label{exercise:ch01-附加良序-2}
(a) 设 $J$ 与 $E$ 都是良序集，设 $h: J \longrightarrow E$。证明以下两个条件等价：
\begin{enumerate}[(i)]
\item $h$ 是保序的并且它的像是 $E$ 或 $E$ 的一个截.
\item $h(\alpha) = [E - h(S_{\alpha})]$ 的最小元素，其中 $\alpha$ 取遍 $J$ 中所有元素.
\end{enumerate}
[提示：证明上述每一条件都蕴涵着 $h(S_{\alpha})$ 为 $E$ 的一个截，因而必定是在 $h(\alpha)$ 处的一个截。]

(b) 设 $E$ 是一个良序集。证明 $E$ 不具有与 $E$ 有相同序型的截，$E$ 的任何两个不同的截也没有相同的序型。[提示：对于任意给定的 $J$，最多只有一个从 $J$ 到 $E$ 中的保序对应的像集为 $E$ 或 $E$ 的一个截。]
\end{exercise}

\begin{exercise}\label{exercise:ch01-附加良序-3}
设 $J$ 和 $E$ 都是良序集，$k: J \longrightarrow E$ 是一个保序对应。应用习题 1 及习题 2 证明：或者 $J$ 与 $E$ 有相同的序型，或者 $J$ 与 $E$ 的一个截有相同的序型。[提示：选取 $e_{0} \in E$。据以下归纳公式定义 $h: J \longrightarrow E$：当 $h(S_{\alpha}) \neq E$ 时，
\begin{align*}
h(\alpha) = [E - h(S_{\alpha})] \text{ 的最小元};
\end{align*}
否则，$h(\alpha) = e_{0}$。证明对于所有 $\alpha$，$h(\alpha) \leqslant k(\alpha)$；从而推得，对于所有 $\alpha$，$h(S_{\alpha}) \neq E$。]
\end{exercise}

\begin{exercise}\label{exercise:ch01-附加良序-4}
根据习题 1$\sim$3 证明以下结论：
\begin{enumerate}[(a)]
\item 设 $A$ 和 $B$ 都是良序集，那么以下的三种可能中恰有其中之一成立：$A$ 与 $B$ 有相同的序型；$A$ 与 $B$ 的一个截有相同的序型；$B$ 与 $A$ 的一个截有相同的序型。[提示：仿照第\cref{section:ch01-良序集}节习题 8 构造一包含 $A$ 和 $B$ 的良序集，然后再应用上一个习题的结论。]
\item 设 $A$ 与 $B$ 都是不可数的良序集，使得 $A$ 和 $B$ 的每一个截都是可数集。证明 $A$ 与 $B$ 有相同的序型.
\end{enumerate}
\end{exercise}

\begin{exercise}\label{exercise:ch01-附加良序-5}
设 $X$ 是一个集合，$\mathcal{A}$ 是所有 $(A, <)$ 偶对的族，其中 $A$ 是 $X$ 的一个子集，$<$ 为 $A$ 上的一个良序。当 $(A, <)$ 是 $(A', <')$ 的一个截时，定义
\begin{align*}
(A, <) \prec (A', <').
\end{align*}
\begin{enumerate}[(a)]
\item 证明 $\prec$ 是 $\mathcal{A}$ 中的一个严格偏序.
\item 设 $\mathcal{B}$ 关于 $\prec$ 是 $\mathcal{A}$ 的一个全序子集。令 $B'$ 为所有使得 $(B, <) \in \mathcal{B}$ 的集合 $B$ 的并。再定义 $<'$ 为所有使得 $(B, <) \in \mathcal{B}$ 的关系 $<$ 之并。证明 $(B', <')$ 是一个良序集.
\end{enumerate}
\end{exercise}

\begin{exercise}\label{exercise:ch01-附加良序-6}
应用习题 1$\sim$5 证明：

\begin{theorem}[极大原理等价于良序定理]\label{theorem:ch01-极大等价良序}
极大原理等价于良序定理。
\end{theorem}
\end{exercise}

\begin{exercise}\label{exercise:ch01-附加良序-7}
应用习题 1$\sim$5 证明：

\begin{theorem}[选择公理等价于良序定理]\label{theorem:ch01-选择等价良序}
\cref{axiom:ch01-选择公理}等价于良序定理。
\end{theorem}

证明：设 $X$ 是一个集合。$c$ 是 $X$ 的所有非空子集的族的一个取定的选择函数。设 $T$ 是 $X$ 的一个子集，$<$ 是 $T$ 中的一个关系。如果 $<$ 是 $T$ 上的一个良序，并且对于每一个 $x \in T$ 有
\begin{align*}
x = c(X - S_{x}(T)),
\end{align*}
则称 $(T, <)$ 是一个塔（tower），其中 $S_{x}(T)$ 为 $T$ 在 $x$ 处的一个截。
\begin{enumerate}[(a)]
\item 设 $(T_{1}, <_{1})$ 与 $(T_{2}, <_{2})$ 是 $X$ 中的两个塔。证明：或者两个有序集相同，或者其中之一为另一个的一个截。[提示：必要时调换下标，总可以假设 $h: T_{1} \longrightarrow T_{2}$ 保序，并且 $h(T_{1})$ 等于 $T_{2}$ 或者 $h(T_{1})$ 是 $T_{2}$ 的一个截。应用习题 2 证明，对于所有 $x$，$h(x) = x$。]
\item 若 $(T, <)$ 是 $X$ 中的一个塔并且 $T \neq X$，证明存在 $X$ 的一个塔使得 $(T, <)$ 为其一个截.
\item 设 $\{(T_{k}, <_{k}) \mid k \in K\}$ 为由 $X$ 中所有塔所组成的族。设
\begin{align*}
T = \bigcup_{k \in K} T_{k} \quad \text{并且} \quad < = \bigcup_{k \in K} (<_{k}).
\end{align*}
证明 $(T, <)$ 是 $X$ 的一个塔。从而得到 $T = X$.
\end{enumerate}
\end{exercise}

\begin{exercise}\label{exercise:ch01-附加良序-8}
应用习题 1$\sim$4 中的结果，按照以下方式构造一个不可数良序集。设 $\mathcal{A}$ 是所有偶对 $(A, <)$ 的族，其中 $A$ 是 $\mathbb{Z}_{+}$ 的一个子集，$<$ 是 $A$ 中的一个良序。（允许 $A$ 为空集。）如果 $(A, <)$ 与 $(A', <')$ 有相同的序型，则定义 $(A, <) \sim (A', <')$。显然，这是一个等价关系。设 $[(A, <)]$ 为 $(A, <)$ 的等价类，$E$ 为所有等价类的族。当 $[(A, <)]$ 与 $[(A', <')]$ 的一个截有相同的序型时，定义
\begin{align*}
[(A, <)] \ll [(A', <')].
\end{align*}
\begin{enumerate}[(a)]
\item 证明关系 $\ll$ 的定义是确切的，并且是 $E$ 中的一个全序。等价类 $[(\varnothing, \varnothing)]$ 是 $E$ 的最小元.
\item 证明若 $\alpha = [(A, <)]$ 为 $E$ 的一个元素，则 $(A, <)$ 与 $E$ 在 $\alpha$ 处的一个截 $S_{\alpha}(E)$ 有相同的序型。[提示：定义一个映射 $f: A \longrightarrow E$，使得对于每一个 $x \in A$，$f(x) = [(S_{x}(A), <\text{的限制})]$。]
\item 证明 $\ll$ 是 $E$ 中的一个良序.
\item 证明 $E$ 是不可数的。[提示：若 $h: E \longrightarrow \mathbb{Z}_{+}$ 为一一对应，则 $h$ 导出 $\mathbb{Z}_{+}$ 上一个良序。]
\end{enumerate}
以任何一个良序集 $X$ 更换 $\mathbb{Z}_{+}$，经相仿的讨论，证明（不用\cref{axiom:ch01-选择公理}）存在一个良序集 $E$，使得它的基数大于 $X$ 的基数。

本习题说明，不用\cref{axiom:ch01-选择公理}我们也能以一个简练的方式构造不可数良序集，从而得到最小不可数良序集。然而，这种构造使得最小不可数良序集本身所应当呈现出的性质有所削弱。我们通常使用的 $S_{\Omega}$ 关键性质是 $S_{\Omega}$ 的每一可数子集在 $S_{\Omega}$ 中有上界。事实上，它依赖于可数集的可数并依然是可数的这一事实。而这一事实的证明（如果你仔细考察）已涉及无穷多次的任意选取，即它依赖于\cref{axiom:ch01-选择公理}。

换言之，不用\cref{axiom:ch01-选择公理}我们也可以构造最小不可数良序集，但却难以用它来处理所有问题。
\end{exercise}

\end{document}